\subsection{A.4 continued: exact scope of the homotopy comparisons}

\paragraph{A-018. Proved: integral homotopy determination.}
For a connected nilpotent space $X$ of finite type, the full weighted
PD cochain algebra over $\mathbb Z$, with its differential, weight
grading and $\Gamma_{\mathbb Z}(T)$-algebra structure, determines
the weak homotopy type of $X$. We regard that strict algebra as an
$E_\infty$-algebra and apply the symmetric monoidal derived shadow
construction of Statement~\ref{sa08:D006E}. It recovers
$C^*(X;\mathbb Z)$ as an $E_\infty$-algebra. Mandell's Main Theorem
in \cite[p.~213]{A_Mandell2006} then gives the asserted
determination. This is the entire reduction needed to apply that
theorem; it does not replace $E_\infty$-equivalence by isomorphism
of cohomology rings. Nor does it assert that the cochain functor is
full. Mandell's Theorems~0.1 and~0.2, pp.~214--215, describe its
retraction and the additional maps in its right adjoint.
\hfill$\square$

\paragraph{A-019. Known, with the comparison to our object specified.}
The following distinctions are needed for the proposed binomial
bridge.

\emph{Horel.}
The functor $X\mapsto\mathbb Z^X$ takes values in cosimplicial
binomial rings: its term in cosimplicial degree $n$ is the product
ring $\mathbb Z^{X_n}$ with coordinatewise binomial operations.
Horel's Theorem~5.4 gives a weak equivalence from a connected
nilpotent finite-type $X$ to the derived binomial realization of
this object. Theorem~6.1 identifies the corresponding unit for
connected finite-type spaces with Bousfield--Kan integral
completion. These statements, in \cite{A_Horel}, are stronger
mapping-space statements than object determination alone. They
do not put a degree filtration on every function $X_n\to\mathbb Z$.
Polynomial degree belongs to a chosen numerical presentation;
identifying its associated graded with our PD cochains requires a
further comparison. A-016 excludes the stated strict simplicial
DGA realization of such a comparison.

\emph{To\"en.}
For a simply connected finite-type $X$, To\"en's Theorem~4.1
identifies the homotopy sheaves of its integral affine homotopy
type with $\pi_i(X)\otimes\mathbb H$, where $\mathbb H$ is the
Hilbert additive group. Corollary~5.1 gives
\[
 \pi_i\bigl((X\otimes\mathbb Z)^{\rm uni}(\mathbb Z)\bigr)
 \cong \pi_i(X)\oplus\prod_p\pi_{i+1}(X)^{\wedge}_p.
\]
Corollaries~5.4--5.5 give a functorial retraction, determination of
objects, and a split injection on mapping spaces
\cite{A_Toen2020}. Thus the integral points themselves generally
have extra homotopy groups. The input is a cosimplicial
commutative ring. We do not identify that object, before applying
any comparison, with the strict differential graded algebra of
PD forms. The independent recovery in A-018 supplies exactly the
integral object-determination conclusion used here.

\emph{Cartan.}
Theorem~1 of \cite[\S2, p.~264]{A_Cartan1976} proves comparison
for an exact cochain resolution whose simplicial terms are
weakly contractible. Theorem~2, p.~266, constructs and uniquely
specifies a linear comparison morphism, using exactness and the
hypothesis $A_p^n=0$ when $n>p$. This is a simplexwise vanishing
hypothesis, not a strict cohomological comparison range.
Theorem~3 in \S5 proves multiplicativity on cohomology under its
flat-cycle hypothesis. Section~7, pp.~269--271, gives the PD
construction, the all-$X$ comparison with coefficient tails,
the fixed-weight endpoint $r\geq n$, the binomial product and
recovery over $\mathbb Z$. We independently reprove those
universal conclusions and establish the arbitrary-base,
derived-tail and completion variants where used. These results
do not supply a strict uniqueness theorem for the category in C.3.

\emph{Numerical spaces and filtered circles.}
Ekedahl's numerical models \cite{A_Ekedahl2001} permit numerical
maps of degree greater than one. Such maps need not preserve the
ordinary degree filtration: the substitution $x\mapsto\binom x2$
already doubles the degree of $x$. A filtration on a numerical
presentation therefore needs additional data or comparison maps.
Hedenlund--Moulinos \cite[\S5.2 and Proposition~5.17]{A_HedenlundMoulinos}
give the precise filtered bar description
\[
 \operatorname{fil}_{\deg}\operatorname{Int}(\mathbb Z)
 \simeq L_0\mathbb Z\otimes^{\mathbf L}_
 {\tau_{\geq\bullet}C^*(S^1;\mathbb Z)}L_0\mathbb Z.
\]
Its associated-graded coefficient algebra is a divided-power
algebra. This theorem concerns the filtered bar construction;
it does not identify numerical K\"ahler forms with PD forms.
The explicit difference in A-021 below is visible even on the
circle. These exact comparisons leave no coefficient-associated-
graded argument serving silently as a comparison of cochain
models.

\paragraph{A-020. Proved: coordinate identification with Kageyama's PD simplex.}
The PD simplex in \cite[\S2.1]{A_Kageyama2022}, and in the section
on differential forms on standard simplices of
\cite{A_Kageyama2024}, is our universal relative PD simplex in
cumulative coordinates. Explicitly, define
\[
 Y_i=\sum_{j=i}^nX_j\quad(0\leq i\leq n),\qquad Y_0=T.
\]
This integral triangular change of basis has determinant one and
therefore induces an isomorphism of free PD algebras, of their
relative PD differential modules, and of their exterior algebras.
For $\alpha:[m]\to[n]$ monotone,
\[
 \alpha^*Y_i=\sum_{\alpha(j)\geq i}X_j
 =\begin{cases}
 Y_{\min\{j:\alpha(j)\geq i\}},&\text{if that set is nonempty},\\
 0,&\text{otherwise}.
 \end{cases}
\]
This is exactly the cumulative-coordinate pullback in the opened
sources. It fixes $Y_0=T$, commutes with the relative differential,
and preserves total weight. Hence it is an isomorphism of the
specified simplicial PD DGAs, not just an isomorphism in each
simplex. Their construction can therefore be used with our
weight convention after this explicit change of coordinates.
We use the binomial PD product determined by their displayed
rational embedding: the multiplication displays at source line
528 in the 2022 paper and line 257 in the 2024 paper omit the
numerator factorials. The 2022 paper explicitly imposes zero odd
squares, so its graded symmetric algebra has the exterior
convention used here. No stronger homotopy-recovery claim is
imported from those sources.\hfill$\square$

\paragraph{A-021. Refuted: an immediate identification with the filtered-circle associated graded.}
Give the positive cohomology generator of the Postnikov associated
graded of $C^*(S^1;\mathbb Z)$ weight one. That graded cohomology
is $\mathbb Z\oplus\mathbb Z[-1](-1)$. Even after extending its
coefficients to $S=\Gamma_{\mathbb Z}(T)$, its degree-one graded
$S$-module is $S[-1]$, whereas A-003 gives $I_1$ for the PD
cochains of the circle. These two graded $S$-modules are not
isomorphic.

\emph{Proof.}
Their weight-one and weight-two pieces are free of rank one.
Multiplication by $T$ from weight one to weight two is multiplication
by $1$ for $S[-1]$ and by $2$ for $I_1$. A graded isomorphism
would act by a unit $\pm1$ in each of those pieces; the resulting
square cannot commute. The four possibilities are checked by
\texttt{A\_model\_comparison.py}. Thus an associated-graded
coefficient calculation alone does not identify these models.
\hfill$\square$
